\section{Introduction}
\label{sec:intro}

Judging whether a research idea is new is difficult, and the growing volume of
published work makes it harder. Novelty is not a monotone function of textual
similarity to any single prior work, and separating a genuine advance from a
restatement is in practice a judgment that draws on familiarity with the field.
We therefore study \emph{evidence-conditioned} novelty assessment: grading a
target idea relative to an explicit set of prior works published before it.
Fixing this set makes the basis of each judgment inspectable and turns the
evidence into a variable that can be controlled.

Much existing work takes a paper as the unit of assessment, from citation-pattern
measures \citep{uzzi2013atypical} to document-level embeddings
\citep{cohan2020specter,ostendorff2022scincl} and to systems that review a
submission as a whole \citep{zhu2025deepreview}. A paper bundles contributions of
uneven novelty: a central mechanism may sit beside incidental variants, and one
score over the bundle does not say which contribution it is about. The decision
of what to pursue is made one idea at a time.

Systems that do operate on ideas assemble the judgment as a language-model
workflow, retrieving related work and then prompting a model to compare against
it \citep{wang2024scimon,dasilva2025graphmind,afzal2026beyond}. Both halves are
fragile. The evidence is selected by the system being evaluated, so a verdict
cannot be attributed to the retrieval or to the comparison, and neither can be
measured or improved on its own. The comparison itself is elicited by a prompt
rather than learned against supervision. \citet{schopf2026rinobench} evaluate
this arrangement against expert judgments and report that although the reasoning
such models produce closely mirrors human rationales, that alignment does not
reliably translate into accurate novelty judgments, which diverge significantly
from the expert gold standard.

Such a workflow also compresses the comparison into a single score, and three
distinctions do not survive that compression. Historical coverage can be
\emph{dispersed}: several papers may jointly account for an idea even when no
single paper closely resembles it, so a comparison against the nearest prior
work misses it. Ideas at equal coverage differ in \emph{what remains}: one may
leave only an implementation detail unexplained while another introduces a new
mechanism. And components that are individually covered may form a
\emph{combination} that no single prior work puts together. Each is a question
about a particular part of the idea and the particular prior work responsible
for it, and a score that discards this correspondence cannot be post-processed
to recover it. Assessing novelty calls instead for a representation that keeps
it: which prior work accounts for which part of the target, and what is left
over.

Studying this setting requires ideas paired with the evidence admissible at the
time they were proposed, which existing novelty benchmarks do not provide. We
therefore construct \textsc{IdeaDelta}, a dataset spanning graph machine
learning, computer vision, and time series, in which every entry is an
idea-length statement paired with
temporally admissible priors and annotated at both the idea and target--prior
levels. The pairing supports the proposed task and makes controlled changes to
the supplied evidence possible.

We propose Evidence-Grounded Residual Decomposition (EGRD), which decomposes a
target against ranked historical evidence into what that evidence accounts for
and what it leaves, and predicts an ordered novelty grade from the two. For
larger evidence sets, a conservative blockwise inference rule extends the
evidence budget without retraining. Our contributions are:
\begin{itemize}
  \item \textbf{\textsc{IdeaDelta}, a multi-domain, evidence-paired dataset.}
  We provide ideas paired
  with temporally admissible priors and annotations of coverage, residual
  contribution, and novelty across three research areas. Because the evidence is
  supplied rather than retrieved, it can be intervened on: permuting the priors
  while holding the target fixed drops the model to near-chance in every domain.
  \item \textbf{Unit-level coverage over dispersed evidence.} We model how much
  of the target the supplied priors account for at the level of latent semantic
  units rather than of the idea as a whole. A learned support matrix aligns each
  unit with each prior, and noisy-OR aggregation lets several partially relevant
  papers jointly cover a unit that none of them covers alone.
  \item \textbf{Typed residuals.} We represent what the priors leave unaccounted
  for as two distinct quantities: unit residuals, for content that no prior
  supports, and interaction residuals, for combinations of individually covered
  components that no single prior puts together. Both are read off the same
  support matrix, so the distinction between an unexplored contribution and an
  effective recombination is produced by the model rather than asserted.
\end{itemize}

Appendix~\ref{app:related} situates EGRD against automated idea generation,
prior novelty judges, and decomposed verification against evidence.
